\textbf{Phần 1. Trắc nghiệm lựa chọn}
\setcounterpageref{ex}{0}

\begin{ex} %[5D10N62]
Muốn tính vận tốc của một chuyển động đều, ta lấy
\choice
{Quãng đường nhân với thời gian}
{\True Quãng đường chia cho thời gian}
{Thời gian chia cho quãng đường}
{Quãng đường cộng với thời gian}
\loigiai{
$v = s : t$.\\
\textbf{Đáp án của câu hỏi là B}
}
\end{ex}

\begin{ex} %[5D10N62]
$3$ giờ $25$ phút $+ 2$ giờ $45$ phút bằng
\choice
{$5$ giờ $70$ phút}
{$6$ giờ $15$ phút}
{\True $6$ giờ $10$ phút}
{$5$ giờ $10$ phút}
\loigiai{
$3$ giờ $25$ phút $+ 2$ giờ $45$ phút $= 5$ giờ $70$ phút $= 6$ giờ $10$ phút.\\
\textbf{Đáp án của câu hỏi là C}
}
\end{ex}

\begin{ex} %[5D10H62]
Một người đi xe máy đi được $105\text{ km}$ trong $2$ giờ $30$ phút. Vận tốc của người đó là
\choice
{$40\text{ km/h}$}
{\True $42\text{ km/h}$}
{$45\text{ km/h}$}
{$50\text{ km/h}$}
\loigiai{
$105 : 2{,}5 = 42\text{ (km/h)}$.\\
\textbf{Đáp án của câu hỏi là B}
}
\end{ex}

\begin{ex} %[5D10H62]
Một ô tô đi với vận tốc $54\text{ km/h}$ trong $1$ giờ $40$ phút. Quãng đường ô tô đi được là
\choice
{$80\text{ km}$}
{$95\text{ km}$}
{$100\text{ km}$}
{\True $90\text{ km}$}
\loigiai{
$1$ giờ $40$ phút $= \dfrac{5}{3}$ giờ.\\
$s = 54 \times \dfrac{5}{3} = 90\text{ (km)}$.\\
\textbf{Đáp án của câu hỏi là D}
}
\end{ex}

\begin{ex} %[5D10H62]
Hai tỉnh A và B cách nhau $140\text{ km}$. Hai ô tô đi ngược chiều nhau cùng khởi hành và gặp nhau sau $1$ giờ $45$ phút. Tổng vận tốc của hai xe là
\choice
{\True $80\text{ km/h}$}
{$70\text{ km/h}$}
{$85\text{ km/h}$}
{$90\text{ km/h}$}
\loigiai{
$1$ giờ $45$ phút $= 1{,}75$ giờ.\\
Tổng vận tốc: $140 : 1{,}75 = 80\text{ (km/h)}$.\\
\textbf{Đáp án của câu hỏi là A}
}
\end{ex}

\begin{ex} %[5D10H62]
Một máy bay bay với vận tốc $850\text{ km/h}$ trên quãng đường $2\,125\text{ km}$. Thời gian máy bay bay hết quãng đường là
\choice
{$2$ giờ $15$ phút}
{$2$ giờ $45$ phút}
{\True $2$ giờ $30$ phút}
{$3$ giờ}
\loigiai{
$t = 2\,125 : 850 = 2{,}5\text{ giờ} = 2$ giờ $30$ phút.\\
\textbf{Đáp án của câu hỏi là C}
}
\end{ex}

\begin{ex} %[5D10V62]
Một người đi xe máy từ A đến B với vận tốc $36\text{ km/h}$. Sau khi xe máy đi được $1$ giờ $15$ phút thì một ô tô cũng khởi hành từ A đuổi theo với vận tốc $54\text{ km/h}$. Thời gian ô tô đuổi kịp xe máy là
\choice
{$2$ giờ}
{$2$ giờ $15$ phút}
{$3$ giờ}
{\True $2$ giờ $30$ phút}
\loigiai{
$1$ giờ $15$ phút $= 1{,}25$ giờ.\\
Khoảng cách xe máy đã đi: $36 \times 1{,}25 = 45\text{ (km)}$.\\
Hiệu vận tốc: $54 - 36 = 18\text{ (km/h)}$.\\
Thời gian đuổi kịp: $45 : 18 = 2{,}5\text{ giờ} = 2$ giờ $30$ phút.\\
\textbf{Đáp án của câu hỏi là D}
}
\end{ex}

\begin{ex} %[5D10V62]
Một ca nô đi xuôi dòng trên khúc sông từ A đến B dài $60\text{ km}$ hết $2$ giờ. Biết vận tốc dòng nước chảy là $4\text{ km/h}$. Thời gian ca nô đi ngược dòng từ B về A là
\choice
{$2{,}5\text{ giờ}$}
{$3$ giờ}
{\True $2$ giờ $43$ phút (hay $\dfrac{30}{11}$ giờ)}
{$2$ giờ $30$ phút}
\loigiai{
Vận tốc xuôi dòng: $60 : 2 = 30\text{ km/h}$.\\
Vận tốc ngược dòng: $30 - 4 \times 2 = 22\text{ km/h}$.\\
Thời gian ngược dòng: $60 : 22 = \dfrac{30}{11}\text{ giờ} \approx 2$ giờ $43$ phút.\\
\textbf{Đáp án của câu hỏi là C}
}
\end{ex}

\begin{ex} %[5D10V62]
Một người đi từ A đến B với vận tốc $40\text{ km/h}$, khi về đi với vận tốc $50\text{ km/h}$. Thời gian về ít hơn thời gian đi là $30$ phút. Quãng đường AB dài là
\choice
{\True $100\text{ km}$}
{$80\text{ km}$}
{$120\text{ km}$}
{$150\text{ km}$}
\loigiai{
Tỉ số vận tốc đi và về là $\dfrac{40}{50} = \dfrac{4}{5} \Rightarrow$ Tỉ số thời gian đi và về là $\dfrac{5}{4}$.\\
Hiệu thời gian là $30$ phút $= 0{,}5$ giờ.\\
Thời gian đi: $0{,}5 : (5 - 4) \times 5 = 2{,}5$ giờ.\\
Quãng đường AB: $40 \times 2{,}5 = 100\text{ (km)}$.\\
\textbf{Đáp án của câu hỏi là A}
}
\end{ex}

\begin{ex} %[5D10V62]
Một con tàu hỏa chạy qua một cây cầu dài $360\text{ m}$ mất $30$ giây và vượt qua một cột điện ven đường mất $12$ giây. Vận tốc của tàu hỏa là
\choice
{$15\text{ m/s}$}
{\True $20\text{ m/s}$}
{$25\text{ m/s}$}
{$30\text{ m/s}$}
\loigiai{
Thời gian tàu đi hết chiều dài cây cầu $360\text{ m}$ là: $30 - 12 = 18$ giây.\\
Vận tốc của tàu: $360 : 18 = 20\text{ (m/s)}$.\\
\textbf{Đáp án của câu hỏi là B}
}
\end{ex}

\begin{ex} %[5D10V62]
Một người đi bộ một vòng quanh hồ với vận tốc $4{,}5\text{ km/h}$ hết $36$ phút. Nếu người đó đi xe đạp với vận tốc $13{,}5\text{ km/h}$ thì đi hết một vòng hồ đó trong bao lâu?
\choice
{$10$ phút}
{$15$ phút}
{$18$ phút}
{\True $12$ phút}
\loigiai{
Vận tốc xe đạp gấp vận tốc đi bộ: $13{,}5 : 4{,}5 = 3$ lần.\\
Thời gian đi xe đạp: $36 : 3 = 12$ phút.\\
\textbf{Đáp án của câu hỏi là D}
}
\end{ex}

\begin{ex} %[5D10C62]
Hai xe cùng xuất phát từ A đến B. Xe thứ nhất đi với vận tốc $45\text{ km/h}$, xe thứ hai đi với vận tốc $60\text{ km/h}$. Xe thứ nhất đến B sau xe thứ hai $40$ phút. Tính quãng đường AB.
\choice
{\True $120\text{ km}$}
{$100\text{ km}$}
{$140\text{ km}$}
{$150\text{ km}$}
\loigiai{
Tỉ số vận tốc xe 1 và xe 2 là $\dfrac{45}{60} = \dfrac{3}{4} \Rightarrow$ Tỉ số thời gian là $\dfrac{4}{3}$.\\
Hiệu thời gian: $40$ phút $= \dfrac{2}{3}$ giờ.\\
Thời gian xe 1 đi: $\dfrac{2}{3} : (4 - 3) \times 4 = \dfrac{8}{3}$ giờ.\\
Quãng đường AB: $45 \times \dfrac{8}{3} = 120\text{ (km)}$.\\
\textbf{Đáp án của câu hỏi là A}
}
\end{ex}

\textbf{Phần 2. Trắc nghiệm đúng sai}
\setcounterpageref{ex}{0}

\begin{ex} %[5D10H62]
Hai địa điểm A và B cách nhau $120\text{ km}$. Xe máy đi từ A với vận tốc $35\text{ km/h}$, ô tô đi từ B với vận tốc $55\text{ km/h}$, hai xe cùng khởi hành lúc $7$ giờ $15$ phút đi ngược chiều nhau. Xét tính đúng sai:
\choiceTF
{\True Sau mỗi giờ hai xe đi được quãng đường $90\text{ km}$.}
{\True Hai xe gặp nhau sau $1$ giờ $20$ phút ($1{,}33$ giờ).}
{Thời điểm hai xe gặp nhau là $8$ giờ $45$ phút.}
{\True Thời điểm hai xe gặp nhau là $8$ giờ $35$ phút.}
\loigiai{
\begin{itemize}
    \item Tổng vận tốc: $35 + 55 = 90\text{ km/h}$. Đúng.
    \item Thời gian gặp: $120 : 90 = \dfrac{4}{3}\text{ giờ} = 1$ giờ $20$ phút. Đúng.
    \item Thời điểm gặp: $7$h $15$p $+ 1$h $20$p $= 8$h $35$p $\neq 8$h $45$p. Ý 3 Sai, Ý 4 Đúng.
\end{itemize}
\begin{itemize}
    \item a) \textbf{Đúng}.
    \item b) \textbf{Đúng}.
    \item c) \textbf{Sai}.
    \item d) \textbf{Đúng}.
\end{itemize}
}
\end{ex}

\begin{ex} %[5D10V62]
Một người đi xe đạp từ A đến B mất $3$ giờ với vận tốc $12\text{ km/h}$. Khi từ B trở về A người đó đi xe máy với vận tốc $36\text{ km/h}$. Xét tính đúng sai:
\choiceTF
{\True Quãng đường AB dài $36\text{ km}$.}
{\True Thời gian đi từ B về A bằng xe máy là $1$ giờ.}
{Tổng thời gian cả đi và về là $4$ giờ $30$ phút.}
{\True Vận tốc trung bình cả đi và về là $18\text{ km/h}$.}
\loigiai{
\begin{itemize}
    \item $s = 12 \times 3 = 36\text{ km}$. Đúng.
    \item Thời gian về: $36 : 36 = 1$ giờ. Đúng.
    \item Tổng thời gian: $3 + 1 = 4$ giờ $\neq 4$h $30$p. Sai.
    \item Vận tốc trung bình: $(36 + 36) : (3 + 1) = 72 : 4 = 18\text{ km/h}$. Đúng.
\end{itemize}
\begin{itemize}
    \item a) \textbf{Đúng}.
    \item b) \textbf{Đúng}.
    \item c) \textbf{Sai}.
    \item d) \textbf{Đúng}.
\end{itemize}
}
\end{ex}

\textbf{Phần 3. Trả lời ngắn}
\setcounterpageref{ex}{0}

\begin{ex} %[5D10H62]
Một người đi bộ với vận tốc $4{,}8\text{ km/h}$ trong $1$ giờ $15$ phút. Tính quãng đường người đó đi được theo ki-lô-mét.
\shortans{6}
\loigiai{
\textbf{Đáp số:} $6$\\
$4{,}8 \times 1{,}25 = 6\text{ km}$.
}
\end{ex}

\begin{ex} %[5D10H62]
Một ô tô đi quãng đường $210\text{ km}$ với vận tốc $60\text{ km/h}$. Tính thời gian ô tô đi theo giờ (dùng dấu phẩy nếu có).
\shortans{3,5}
\loigiai{
\textbf{Đáp số:} $3{,}5$\\
$210 : 60 = 3{,}5\text{ giờ}$.
}
\end{ex}

\begin{ex} %[5D10V62]
Hai xe máy đi ngược chiều từ hai tỉnh cách nhau $135\text{ km}$ cùng khởi hành lúc $8$ giờ với vận tốc $42\text{ km/h}$ và $48\text{ km/h}$. Hỏi sau bao nhiêu phút hai xe gặp nhau?
\shortans{90}
\loigiai{
\textbf{Đáp số:} $90$\\
$135 : (42 + 48) = 1{,}5\text{ giờ} = 90\text{ phút}$.
}
\end{ex}

\begin{ex} %[5D10V62]
Một người đi xe máy từ A đuổi theo một người đi xe đạp khởi hành cùng lúc cách A $24\text{ km}$. Vận tốc xe máy là $40\text{ km/h}$, xe đạp là $16\text{ km/h}$. Hỏi sau bao nhiêu giờ xe máy đuổi kịp xe đạp?
\shortans{1}
\loigiai{
\textbf{Đáp số:} $1$\\
$24 : (40 - 16) = 24 : 24 = 1\text{ giờ}$.
}
\end{ex}

\begin{ex} %[5D10V62]
Một con chim bồ câu đưa thư bay quãng đường $120\text{ km}$ với vận tốc $50\text{ km/h}$. Tính thời gian chim bồ câu bay theo phút.
\shortans{144}
\loigiai{
\textbf{Đáp số:} $144$\\
$120 : 50 = 2{,}4\text{ giờ} = 144\text{ phút}$.
}
\end{ex}

\begin{ex} %[5D10C62]
Một chiếc thuyền đi xuôi dòng một khúc sông dài $40\text{ km}$ mất $1$ giờ $15$ phút. Vận tốc dòng nước là $4\text{ km/h}$. Hỏi thuyền đi ngược dòng khúc sông đó mất bao nhiêu giờ (dùng dấu phẩy nếu có)?
\shortans{1,67 hoặc \dfrac{5}{3}}
\loigiai{
\textbf{Đáp số:} $1{,}67$ hoặc $\dfrac{5}{3}$\\
Vận tốc xuôi dòng: $40 : 1{,}25 = 32\text{ km/h}$.\\
Vận tốc ngược dòng: $32 - 4 \times 2 = 24\text{ km/h}$.\\
Thời gian ngược dòng: $40 : 24 = \dfrac{5}{3}\text{ giờ} = 1\text{ giờ } 40\text{ phút} \approx 1{,}67\text{ giờ}$.
}
\end{ex}
